\subsection{The transcendence degree of an extension}

THEOREM 3. --- Let E be an extension of a field K. All transcendence bases of E over K have the same cardinal.

It is enough to prove the inequality \(\operatorname{Card}(\mathbf{B}) \geqslant \operatorname{Card}(\mathbf{B}')\) , when \(\mathbf{B}\) and \(\mathbf{B}'\) are two transcendence bases of \(\mathbf{E}\) over \(\mathbf{K}\) , and we may suppose that \(\mathbf{B}'\) is not empty. We first suppose that \(\mathbf{B}\) is finite and use induction on its cardinal \(n\) ; for \(n = 0\) , \(\mathbf{E}\) is algebraic over \(\mathbf{K}\) and \(\mathbf{B}'\) is empty. Suppose now that \(n \geqslant 1\) ; given \(x \in \mathbf{B}'\) , the exchange theorem provides a subset \(\mathbf{C}\) of \(\mathbf{B}\) such that \(x \notin \mathbf{C}\) and \(\{x\} \cup \mathbf{C}\) is a transcendence basis of \(\mathbf{E}\) over \(\mathbf{K}\) ; we have \(\mathbf{C} \neq \mathbf{B}\) by Prop. 7 whence \(\operatorname{Card} \mathbf{C} < n\) . Put \(\mathbf{K}_1 = \mathbf{K}(x)\) and \(\mathbf{C}' = \mathbf{B}' - \{x\}\) ; then \(\mathbf{C}\) and \(\mathbf{C}'\) are algebraically free over the field \(\mathbf{K}_1\) (V, p.~107, Prop. 4) and \(\mathbf{E}\) is algebraic both over \(\mathbf{K}_1(\mathbf{C}) = \mathbf{K}(\mathbf{C} \cup \{x\})\) and \(\mathbf{K}_1(\mathbf{C}') = \mathbf{K}(\mathbf{B}')\) . In other words, \(\mathbf{C}\) and \(\mathbf{C}'\) are two transcendence bases of \(\mathbf{E}\) over \(\mathbf{K}_1\) . Since \(\operatorname{Card}(\mathbf{C}) < n\) , the induction hypothesis implies the inequality \(\operatorname{Card}(\mathbf{C}') \leqslant \operatorname{Card}(\mathbf{C}) \leqslant n - 1\) , whence \(\operatorname{Card}(\mathbf{B}') \leqslant n = \operatorname{Card}(\mathbf{B})\) .

Suppose now that B is infinite. Every \(x \in B\) is algebraic over \(\mathrm{K}(\mathrm{B}')\) and so there exists a finite subset \(\mathrm{S}(x)\) of \(B'\) such that x is algebraic over \(\mathrm{K}(\mathrm{S}(x))\) . Write \(\mathrm{S} = \bigcup_{x \in \mathrm{B}} \mathrm{S}(x)\) , then \(S \subset B'\) and since B is infinite, we have \(\operatorname{Card}(\mathrm{S}) \leqslant \operatorname{Card}(\mathrm{B})\)

(E, III, p.~49, Cor. 3). But since every element of B is algebraic over K(S) and E is algebraic over K(B), we conclude that E is algebraic over K(S) (V, p.~19, Prop. 3). Now Prop. 7 implies that S = B', whence the desired inequality Card(B') ≤ Card(B).

DEFINITION 4. --- Let E be an extension of a field K. The cardinal of every transcendence basis of E over K is called the transcendence degree of E over K and written tr. deg \(_{K}\) E.

Th. 2 and 3 and Def. 4 imply the following corollaries, where E denotes an extension of a field K, of finite transcendence degree n.

COROLLARY 1. --- Let S be a subset of E such that E is algebraic over K(S). Then Card(S) ≥ n ; if the cardinal of S is equal to n, then S is algebraically free over K (hence it is then a transcendence basis of E over K).

COROLLARY 2. --- Suppose that \(\mathrm{E} = \mathrm{K}(x_1, \ldots, x_m)\) , then \(m \geqslant n\) ; if moreover \(m = n\) , then \((x_1, \ldots, x_m)\) is a pure basis of \(\mathrm{E}\) over \(\mathrm{K}\) , and \(\mathrm{E}\) is then a pure extension of \(\mathrm{K}\) .

COROLLARY 3. --- Every subset of E which is algebraically free over K has at most n elements, and if it has exactly n elements, it is a transcendence basis of E over K.

THEOREM 4. --- Let K, E and F be three fields such that \(K \subset E \subset F\) . If S is a transcendence basis of E over K and T a transcendence basis of F over E, then \(S \cap T\) is empty and \(S \cup T\) is a transcendence basis of F over K.

For F is algebraic over \(E(T)\) ; moreover, \(E(T)\) is algebraic over the field \(K(S \cup T) = K(S)(T)\) , because E is algebraic over \(K(S)\) (V, p.~18, Cor. 2); therefore (V, p.~19, Prop. 3) F is algebraic over \(K(S \cup T)\) . On the other hand, T being algebraically free over E is so a fortiori over \(K(S)\) , hence (V, p.~107, Prop. 4) \(S \cup T\) is algebraically free over K and \(S \cap T = \varnothing\) .

COROLLARY. --- Let K, E and F be three fields such that \(K \subset E \subset F\) . Then we have

\[
\operatorname{tr} \cdot \deg_ {K} F = \operatorname{tr} \cdot \deg_ {K} E + \operatorname{tr} \cdot \deg_ {E} F.\tag{1}
\]

